\documentclass[11pt]{article}
\usepackage{amsmath,amssymb,amsthm}
\usepackage{algorithm,algorithmic}
\usepackage{hyperref}

\theoremstyle{definition}
\newtheorem{definition}{Definition}
\newtheorem{theorem}{Theorem}
\newtheorem{proposition}{Proposition}
\newtheorem{remark}{Remark}

\DeclareMathOperator{\argmax}{argmax}
\DeclareMathOperator{\argmin}{argmin}

\title{Likelihood Optimization for CCP Reconciliation:\\Mathematical Analysis and PyTorch Implementation Strategy}
\author{Analysis Document}
\date{\today}

\begin{document}
\maketitle

\section{Problem Formulation}

\subsection{Fixed-Point Structure}

We consider the reconciliation likelihood optimization problem where both the extinction probabilities $E(\theta)$ and the reconciliation probabilities $\Pi(\theta)$ are solutions to fixed-point equations parameterized by $\theta = (\delta, \tau, \lambda) \in \mathbb{R}^3_+$.

The fixed-point structure can be written as:
\begin{align}
    E &= F_E(E; \theta) \\
    \Pi &= F_\Pi(\Pi, E; \theta)
\end{align}

Note the key structural property: $E$ does not depend on $\Pi$, but $\Pi$ depends on $E$. This creates a hierarchical fixed-point system.

\subsection{Likelihood Function}

The log-likelihood function is defined as:
\begin{equation}
    \mathcal{L}(\theta) = \log\left(\sum_{s \in \mathcal{S}} \Pi^*_{\text{root},s}(\theta)\right)
\end{equation}
where $\Pi^*(\theta)$ denotes the fixed-point solution and $\mathcal{S}$ is the set of species branches.

Alternatively, following the ALE formulation:
\begin{equation}
    \mathcal{L}(\theta) = \log\left(\frac{\sum_{s \in \mathcal{S}} p_s^O \Pi^*_{\text{root},s}(\theta)}{\sum_{s \in \mathcal{S}} p_s^O (1 - E^*_s(\theta))}\right)
\end{equation}
where $p_s^O$ is the origination probability on branch $s$.

\subsection{Key Observation}

The likelihood depends only on:
\begin{itemize}
    \item A single row of $\Pi$: $\Pi_{\text{root},:}$
    \item The full vector $E$ (if using the normalized likelihood)
\end{itemize}

This sparsity structure is crucial for efficient gradient computation.

\section{Gradient Computation via Implicit Differentiation}

\subsection{Implicit Function Theorem}

At the fixed point, we have:
\begin{align}
    E^* - F_E(E^*; \theta) &= 0 \\
    \Pi^* - F_\Pi(\Pi^*, E^*; \theta) &= 0
\end{align}

By the implicit function theorem, the derivatives of the fixed points with respect to $\theta$ satisfy:
\begin{align}
    \frac{\partial E^*}{\partial \theta} &= (I - \frac{\partial F_E}{\partial E})^{-1} \frac{\partial F_E}{\partial \theta} \\
    \frac{\partial \Pi^*}{\partial \theta} &= (I - \frac{\partial F_\Pi}{\partial \Pi})^{-1} \left(\frac{\partial F_\Pi}{\partial \theta} + \frac{\partial F_\Pi}{\partial E} \frac{\partial E^*}{\partial \theta}\right)
\end{align}

\subsection{Adjoint Method for Efficient Computation}

Since we only need $\frac{\partial \mathcal{L}}{\partial \theta}$, we can use the adjoint method to avoid computing the full Jacobians.

Define the adjoint variables:
\begin{align}
    v_\Pi &= \frac{\partial \mathcal{L}}{\partial \Pi^*} \\
    v_E &= \frac{\partial \mathcal{L}}{\partial E^*} + v_\Pi^\top \frac{\partial F_\Pi}{\partial E}
\end{align}

Then the gradient is:
\begin{equation}
    \frac{\partial \mathcal{L}}{\partial \theta} = v_E^\top \frac{\partial F_E}{\partial \theta} + v_\Pi^\top \frac{\partial F_\Pi}{\partial \theta}
\end{equation}

The adjoint variables satisfy the linear systems:
\begin{align}
    (I - \frac{\partial F_\Pi}{\partial \Pi})^\top v_\Pi &= \frac{\partial \mathcal{L}}{\partial \Pi^*}^\top \\
    (I - \frac{\partial F_E}{\partial E})^\top v_E &= \frac{\partial \mathcal{L}}{\partial E^*}^\top + \left(\frac{\partial F_\Pi}{\partial E}\right)^\top v_\Pi
\end{align}

\section{PyTorch Implementation Strategy}

\subsection{Custom Autograd Function}

We implement a custom PyTorch autograd function that:
\begin{enumerate}
    \item \textbf{Forward pass}: Computes the fixed points $E^*$, $\Pi^*$ and the likelihood $\mathcal{L}$
    \item \textbf{Backward pass}: Solves the adjoint equations to compute $\nabla_\theta \mathcal{L}$
\end{enumerate}

\begin{algorithm}
\caption{FixedPointLikelihood.forward}
\begin{algorithmic}[1]
\STATE \textbf{Input:} $\theta = (\delta, \tau, \lambda)$
\STATE Compute event probabilities from $\theta$
\STATE Solve $E^* = F_E(E^*; \theta)$ via fixed-point iteration
\STATE Solve $\Pi^* = F_\Pi(\Pi^*, E^*; \theta)$ via fixed-point iteration
\STATE Compute $\mathcal{L} = \log\sum_s \Pi^*_{\text{root},s}$
\STATE Save $E^*$, $\Pi^*$, and intermediate values for backward pass
\STATE \textbf{Return:} $\mathcal{L}$
\end{algorithmic}
\end{algorithm}

\begin{algorithm}
\caption{FixedPointLikelihood.backward}
\begin{algorithmic}[1]
\STATE \textbf{Input:} $\text{grad\_output}$ (scalar)
\STATE Compute $\frac{\partial \mathcal{L}}{\partial \Pi^*}$ (sparse, only root row)
\STATE Solve adjoint equation for $v_\Pi$ using conjugate gradient
\STATE Compute $v_E$ from $v_\Pi$
\STATE Compute $\nabla_\theta \mathcal{L} = v_E^\top \frac{\partial F_E}{\partial \theta} + v_\Pi^\top \frac{\partial F_\Pi}{\partial \theta}$
\STATE \textbf{Return:} $\nabla_\theta \mathcal{L} \cdot \text{grad\_output}$
\end{algorithmic}
\end{algorithm}

\subsection{Handling Log-Space and Masking}

The key challenge is handling $-\infty$ values in log-space computation:

\begin{enumerate}
    \item \textbf{Masking Strategy}: Track finite positions throughout computation
    \item \textbf{Sparse Gradients}: Only compute gradients for positions that affect the likelihood
    \item \textbf{Stable Operations}: Use logsumexp and careful numerical handling
\end{enumerate}

\subsection{Optimization Algorithm Choices}

Based on the AleRax analysis, we have several options:

\subsubsection{L-BFGS-B}
\begin{itemize}
    \item Quasi-Newton method with box constraints
    \item Good for smooth problems with moderate dimensions
    \item Requires exact gradients
\end{itemize}

\subsubsection{Adam with Softplus Transform}
\begin{itemize}
    \item Transform: $\theta = \text{softplus}(\log\theta)$ ensures positivity
    \item Adaptive learning rates handle different parameter scales
    \item Works well with noisy gradients
\end{itemize}

\subsubsection{Newton's Method}
\begin{itemize}
    \item Requires Hessian computation
    \item Can use finite differences of gradients for Hessian
    \item Faster convergence near optimum
\end{itemize}

\section{Implementation Recommendations}

\subsection{Core Components}

\begin{enumerate}
    \item \textbf{FixedPointCCPLogFunction}: Custom autograd function implementing forward/backward
    \item \textbf{AdjointSolver}: Efficient solver for the adjoint linear systems
    \item \textbf{CCPLogOptimizer}: High-level optimizer class wrapping everything
\end{enumerate}

\subsection{Numerical Considerations}

\begin{enumerate}
    \item \textbf{Convergence Criteria}: Use relative tolerance for fixed-point iterations
    \item \textbf{Gradient Clipping}: Prevent exploding gradients in early iterations
    \item \textbf{Parameter Bounds}: Enforce $\theta > 0$ via transformations
    \item \textbf{Checkpointing}: Save intermediate states for debugging
\end{enumerate}

\subsection{Efficient Computation}

\begin{enumerate}
    \item \textbf{Sparse Operations}: Exploit sparsity in $\frac{\partial \mathcal{L}}{\partial \Pi^*}$
    \item \textbf{Reuse Computations}: Cache matrix-vector products
    \item \textbf{GPU Acceleration}: Use PyTorch's GPU support throughout
    \item \textbf{Batching}: Process multiple parameter sets simultaneously if needed
\end{enumerate}

\section{Theoretical Guarantees}

\subsection{Convergence of Fixed-Point Iterations}

By the Banach fixed-point theorem, the iterations converge if:
\begin{equation}
    \|F'(x)\| < 1
\end{equation}

For the reconciliation problem, this holds when event probabilities are not too extreme.

\subsection{Gradient Accuracy}

The implicit differentiation provides exact gradients (up to numerical precision) when:
\begin{enumerate}
    \item Fixed points are computed to sufficient accuracy
    \item Adjoint equations are solved accurately
    \item No numerical overflow/underflow occurs
\end{enumerate}

\subsection{Local Convergence of Optimization}

Near a local minimum $\theta^*$ with positive definite Hessian:
\begin{itemize}
    \item Gradient descent: Linear convergence with rate $O((1-\mu/L)^k)$
    \item Newton's method: Quadratic convergence
    \item L-BFGS: Superlinear convergence
\end{itemize}

where $\mu$ is the smallest eigenvalue and $L$ is the largest eigenvalue of the Hessian.

\section{Conclusion}

The fixed-point structure of the reconciliation problem admits efficient gradient computation via implicit differentiation and the adjoint method. The key insights are:

\begin{enumerate}
    \item The hierarchical structure ($E$ then $\Pi$) simplifies the adjoint computation
    \item The sparsity of $\frac{\partial \mathcal{L}}{\partial \Pi^*}$ makes the method efficient
    \item Custom PyTorch autograd functions can handle the $-\infty$ masking properly
    \item Modern optimization algorithms (Adam, L-BFGS-B) work well with exact gradients
\end{enumerate}

This approach provides a mathematically principled and computationally efficient solution to the likelihood optimization problem for CCP reconciliation.

\end{document}